% ============================================================================
%  Saussure in HLLSet Algebra --- Part IV: Geographical Linguistics
%  A computational exegesis of Ferdinand de Saussure's
%  Course in General Linguistics, Part Four, in the Boolean lattice
%  and GF(2) ring of HLLSet Algebra
% ----------------------------------------------------------------------------
%  Series:  Saussure in HLLSet Algebra (reverse direction of influence)
%  Build:   pdflatex saussure_hllset_part4_geographical_linguistics.tex
%           (run twice)
%  Notation follows Part I:
%    saussure_hllset_part1_general_principles.tex
%  Part II:
%    saussure_hllset_part2_synchronic_linguistics.tex
%  Part III:
%    saussure_hllset_part3_diachronic_linguistics.tex
%  and the EWM PUB series.
% ============================================================================
\documentclass[11pt,a4paper]{article}

\usepackage[utf8]{inputenc}
\usepackage[T1]{fontenc}
\usepackage{amsmath,amssymb,amsthm}
\usepackage{amscd}
\usepackage{booktabs}
\usepackage{array}
\usepackage{placeins}
\usepackage{hyperref}
\usepackage[margin=1in]{geometry}
\usepackage{enumitem}
\usepackage{xcolor}
\definecolor{hlblue}{RGB}{0,80,160}
\hypersetup{colorlinks=true,linkcolor=hlblue,urlcolor=hlblue,citecolor=hlblue}

% ---- theorems (aligned with the PUB series and Parts I--III) ----------------
\newtheorem{principle}{Principle}[section]
\newtheorem{theorem}[principle]{Theorem}
\newtheorem{definition}[principle]{Definition}
\newtheorem{proposition}[principle]{Proposition}
\newtheorem{corollary}[principle]{Corollary}
\newtheorem{remark}[principle]{Remark}
\theoremstyle{remark}
\newtheorem{observation}[principle]{Observation}
\newtheorem{reading}[principle]{Reading}

% ---- shorthands (aligned with Parts I--III and the PUB series) --------------
\newcommand{\Plane}{\ensuremath{\mathbb{B}}}            % bit plane
\newcommand{\Lat}{\ensuremath{\mathcal{B}}}              % HLLSet (Boolean) lattice
\newcommand{\Tok}{\ensuremath{\mathcal{T}}}              % token space
\newcommand{\HLL}{\ensuremath{H}}                        % an HLLSet
\newcommand{\xor}{\ensuremath{\Delta}}                   % ring addition
\newcommand{\BSS}{\ensuremath{\mathsf{BSS}}}             % BSS measure
\newcommand{\res}{\ensuremath{\mathrm{res}}}             % residual
\newcommand{\spanof}[1]{\ensuremath{\langle #1\rangle}}  % GF(2) span
\newcommand{\supp}{\ensuremath{\mathrm{supp}}}           % support
\newcommand{\Ctx}{\ensuremath{\mathcal{C}}}              % token context
\newcommand{\Ing}{\ensuremath{\mu}}                      % ingestion morphism
\newcommand{\Mat}{\ensuremath{\rho}}                     % materialization morphism
\newcommand{\Lang}{\ensuremath{\mathcal{L}}}             % langue (the lattice system)
\newcommand{\Par}{\ensuremath{\mathcal{P}}}              % parole (usage)
\newcommand{\Sign}{\ensuremath{\mathcal{S}}}             % the sign relation
\newcommand{\Cov}{\ensuremath{\mathrm{cov}}}             % coverage
\newcommand{\Disp}{\ensuremath{\mathcal{D}}}             % the dispatcher
\newcommand{\Sync}{\ensuremath{\mathsf{Syn}}}            % synchronic fiber
\newcommand{\Dia}{\ensuremath{\mathsf{Dia}}}             % diachronic trajectory
\newcommand{\Gen}{\ensuremath{\gamma}}                   % basis generation
\newcommand{\cover}{\ensuremath{\mathrm{cover}}}         % frame cover
\newcommand{\Val}{\ensuremath{\mathcal{V}}}              % value (differential)
\newcommand{\Synt}{\ensuremath{\mathsf{Synt}}}           % syntagmatic reading
\newcommand{\Ass}{\ensuremath{\mathsf{Ass}}}             % associative reading
\newcommand{\Dep}{\ensuremath{\mathsf{Dep}}}             % departed
\newcommand{\Ret}{\ensuremath{\mathsf{Ret}}}             % retained
\newcommand{\New}{\ensuremath{\mathsf{New}}}             % new
\newcommand{\Trans}{\ensuremath{\mathsf{Tr}}}            % recorded transition
\newcommand{\Piv}{\ensuremath{\mathsf{Piv}}}             % pivot
\newcommand{\Ana}{\ensuremath{\mathsf{Ana}}}             % analogical fourth
\newcommand{\Change}{\ensuremath{\mathbf{M}}}            % change-of-basis matrix
% ---- new shorthands for Part IV --------------------------------------------
\newcommand{\World}{\ensuremath{\mathsf{W}}}             % a world
\newcommand{\Iota}{\ensuremath{\iota}}                   % interpretation map
\newcommand{\Bridge}{\ensuremath{\Phi}}                  % bridge between worlds
\newcommand{\Wset}{\ensuremath{\Omega}}                  % family of worlds
\newcommand{\Iso}{\ensuremath{\mathsf{Iso}}}             % isogloss

% ---- title -----------------------------------------------------------------
\title{\textbf{Saussure in HLLSet Algebra}\\[4pt]
\large Part IV: Geographical Linguistics\\[4pt]
\large A Computational Exegesis of the \textit{Course in General Linguistics},\\
Part Four, in the Boolean Lattice and GF(2) Ring of HLLSet Algebra\\[6pt]
\normalsize Fourth paper in the series \textit{Saussure in HLLSet Algebra}}

\author{Alex Mylnikov \and Alexandr Mylnikov Jr. \and Deependra Kumar}
\date{\today}

\begin{document}
\maketitle

\begin{abstract}
Part IV of the series re-founds Part Four of the \textit{Course},
\emph{Geographical Linguistics}, as the multi-world theory of one lattice.
Parts I--III kept one question out of sight: \emph{where} the changes happen.
Saussure's Part Four is that question, and the algebra answers it with the
world layer of the EWM programme. The diversity of languages is the family of
\textbf{worlds} over the one shared plane: hash configurations and codebooks
are dialects, the BSS threshold and the LUT are the boundary of a community,
and one lattice carries all of them at the cost of one LUT per world, not one
lattice per world. The complexity of geographical diversity---dialect
continua, overlapping isoglosses---is the graded geometry of the soft key:
worlds share nodes exactly but read them differently, and an isogloss is a
threshold on a coordinate profile, not a cut. The causes of diversity are
re-founded as Saussure re-founded them: \emph{time} is the primary cause, and
the algebra makes time exact---two communities with a common ancestor frame
diverge by the product of their transport matrices once their pivots stop
being shared; substrata are older generations still readable in the commit
DAG. The propagation of linguistic waves is the \textbf{bridge}
$\Bridge_{w_1\to w_2} = \text{ingest} \circ \text{transcribe} \circ
\text{materialize}$: deterministic, idempotent on builders, monotone, carrying
encodings, not content; a loan is a node whose provenance crosses two worlds,
an inheritance is a node whose provenance stays on one DAG path. The
dispatcher closes the geography: it is the speaker who chooses which world
actualizes the shared langue. The paper ends with the complete dictionary of
Part Four and the residual handed to Part V: reconstruction as time travel
through old frames.
\end{abstract}

\noindent\textbf{Keywords:} Saussure; geographical linguistics; diversity of
languages; dialect; isogloss; substratum; linguistic wave; bridge; world;
hash configuration; codebook; dispatcher; HLLSet; Boolean lattice; GF(2)
ring; computational exegesis.

% ============================================================================
\section{Introduction}
\label{sec:intro}
% ============================================================================

\subsection{Where Part IV stands in the series}

Part I \cite{part1} re-founded the general principles; Part II \cite{part2}
re-founded synchronic linguistics as the differential theory of value; Part
III \cite{part3} released time and re-founded diachronic linguistics as the
theory of the pivot---sound change as extension and rotation, analogy as the
fourth proportional of GF(2), agglutination as the promotion of a view to an
original. All three parts assumed, without saying so, that there is \emph{one}
frame, \emph{one} community, \emph{one} language. Part Four of the
\textit{Course} removes that assumption.

Saussure's Part Four has four chapters: (I) on the diversity of languages;
(II) the complexity of geographical diversity; (III) the causes of
geographical diversity; (IV) the propagation of linguistic waves
\cite[Part IV]{harris1983}. Part IV of this series follows the same order,
after a preliminary section that fixes the world layer.

\subsection{Method, unchanged}

The method is that of Parts I--III \cite{part1,part2,part3}: no metaphor,
principles as theorems or invariants, ambiguity recorded as data. The
translation note of Part I applies unchanged; references to the
\textit{Course} are by part, chapter, and section.

\subsection{Roadmap}

Section~\ref{sec:prelim} fixes the standing definitions and adds the world
layer (worlds, interpretations, bridges). Section~\ref{sec:ch1} re-founds the
diversity of languages as a family of worlds over one plane. Section
\ref{sec:ch2} re-founds the complexity of geographical diversity as graded
overlap and thresholds. Section~\ref{sec:ch3} re-founds the causes of
diversity: time as pivots, separation as unsynchronized generations,
substrata as old frames. Section~\ref{sec:ch4} re-founds the propagation of
linguistic waves as bridges, with loans separated from inheritances by
provenance. Section~\ref{sec:dictionary} collects the dictionary; Section
\ref{sec:adds} states what Part IV adds; Section~\ref{sec:conclusion} hands
the retrospective residual to Part V.

% ============================================================================
\section{Preliminaries: the synchronic and diachronic inheritance, and the world layer}
\label{sec:prelim}
% ============================================================================

We restate the standing definitions with their labels; proofs and motivation
are in \cite{part1,part2,part3}.

\begin{definition}[Plane, lattice, ring, IICA, frames, fibers, pivots]
\label{def:standing}
The \textbf{plane} is $\Plane=\{0,\dots,32767\}$; the \textbf{lattice} is
$\Lat=\mathcal{P}(\Plane)$ under $\cup,\cap,\subseteq$ and the Boolean ring
$\xor=\Delta$, $\cap$. $\Ing:\Tok\to\Lat$ maps each token to its
\textbf{bit-set} and is an \textbf{IICA morphism}
(idempotent, immutable, content-addressed); $\Mat:\Lat\dashrightarrow\Ctx$
materializes through the door, reading each set bit through the LUT. (The
LUT, multi-hashing, and the global gates $G_1,G_2,G_3$ are introduced in
Part I, Section~2.) A
\textbf{basis frame} $F=(B_1,\dots,B_k)$ is
a row-echelon GF(2) basis of observed nodes with distinct leading bits,
stamped with a generation $\Gen$; the hard key of $X\in\spanof{F}$ is
$c\in\mathrm{GF}(2)^k$ with $X=\bigoplus_{i:c_i=1}B_i$, the soft key is
$w_i=\BSS_{B_i}(X)$, and the spill is $s=|X\setminus\cover(F)|$. The
\textbf{synchronic fiber} is $\Sync(t)=(\Lat,F(t),\Gen(t))$; a
\textbf{pivot} is the extension/rotation of Definition~1 of \cite{part3};
the \textbf{transport} of a hard key across a pivot is $c' = c\Change$.
\end{definition}

\subsection{The world layer}

Part IV adds the objects that the EWM programme built for multi-user
representation \cite{companion2026,langue2026}, now read as Saussure's
geography.

\begin{definition}[World; codebook; community]
\label{def:world}
A \textbf{world} is a triple $\World = (h, s, \tau)$: a hash configuration
$(h,s)$ inducing the measurement morphism $\Ing_\World : 2^\Tok \to 2^\Plane$;
a BSS threshold $\tau$ selecting the thresholded node set
$\mathcal{N}_\tau(\World)$; and an \textbf{interpretation}
$\Iota_\World : 2^\Plane \to 2^\Tok$ that reads a node's active bits back as
tokens through the world's LUT \cite{companion2026}. A \textbf{codebook} is
the vocabulary-aware encoder that quantizes a host's encodings to tokens
before $\Ing_\World$ \cite{state2026}. A \textbf{community} is a set of hosts
sharing one world and one generation stamp $\Gen$ of a frame: they speak the
same dialect when their LUTs agree and their frames are synchronized.
\end{definition}

\begin{proposition}[Shared representation, separate interpretation]
\label{prop:sharing}
Two worlds may share the very same node $X \in \Lat$ and yet be distinct:
they are distinct exactly when their interpretations differ on some shared
node,
\[
  \World \neq \World'
  \quad\Longleftrightarrow\quad
  \exists\, X \in \mathcal{N}_\tau(\World)\cap\mathcal{N}_{\tau'}(\World'):
  \Iota_\World(X) \neq \Iota_{\World'}(X) .
\]
All worlds live on the one lattice $\Lat$; what separates them is the LUT,
never the bits.
\end{proposition}

\begin{proof}
This is the representation/interpretation separation of \cite{companion2026}:
the bit vector is semantically inert, and all meaning lives in the LUT, which
is a different object for every hash configuration.
\end{proof}

\begin{proposition}[Holographic multi-worlds]
\label{prop:holographic}
Let $\Wset$ be a family of worlds. Then (i) \emph{common representation:}
every world is read through the same lattice $\Lat$, and the representation
is not duplicated per world; (ii) \emph{separate interpretation:} each world
adds only its LUT and threshold; (iii) \emph{joint closure:} every Boolean
expression over nodes drawn from one world stays in that world's subalgebra.
The cost of hosting $|\Wset|$ worlds is $|\Wset|$ LUTs, not $|\Wset|$
lattices.
\end{proposition}

\begin{proof}
Parts (i)--(iii) are Theorem 2 and the holographic principle of
\cite{companion2026}.
\end{proof}

\begin{definition}[Bridge]
\label{def:bridge}
The \textbf{bridge} from world $\World_1$ to world $\World_2$ is
\[
  \Bridge_{\World_1 \to \World_2}
  \;=\;
  \text{ingest}_{\World_2} \circ \text{transcribe} \circ
  \text{materialize}_{\World_1},
\]
materializing in the source world, transcribing the recovered encodings, and
ingesting under the target world's configuration
\cite{companion2026}.
\end{definition}

\begin{proposition}[Bridge laws]
\label{prop:bridge-laws}
$\Bridge_{\World_1\to\World_2}$ is deterministic (IICA); idempotent on
builders; monotone: $A \subseteq A'$ implies
$\Bridge(A) \subseteq \Bridge(A')$; and in general a strict coarsening:
collisions at materialization can only add atoms to the target. The bridge
carries encodings, not content.
\end{proposition}

\begin{proof}
See \cite{companion2026}: materialize, transcribe, and ingest are
deterministic; composition preserves IICA; materialization is exact on
builders and monotone in the inclusion order.
\end{proof}

% ============================================================================
\section{Chapter I: On the diversity of languages}
\label{sec:ch1}
% ============================================================================

Saussure opens Part Four with the plain fact of diversity: languages differ
from one another in space, as they differ in time \cite[Part IV,
ch.~I]{harris1983}. The diversity is the datum to be explained, not a
deviation from a norm.

\begin{reading}[Diversity is the family of worlds]
\label{read:diversity}
One plane $\Plane$ carries a family $\Wset$ of worlds. A language is a world
together with the frames its community has actually computed: the hash
configuration and codebook are its phonology in the broadest sense, the LUT
is its lexicon, and the frame is its grammar. Two languages are two worlds
over the same plane---two readings of one lattice, in the sense of
Proposition~\ref{prop:sharing}.
\end{reading}

\begin{observation}[Diversity is computed, not given]
\label{obs:diversity}
The lattice does not come divided into languages. What divides it is
\emph{use}: which tokens a community ingests, under which configuration, and
which generations it shares. Diversity is therefore not a property of the
plane but of the family of worlds and their histories. This is Saussure's
point that a language is a social fact: the social fact is the shared LUT and
the synchronized generation.
\end{observation}

\begin{proposition}[One lattice suffices for all languages]
\label{prop:one-lattice}
For any finite or countable family of worlds $\Wset$, the whole family is a
subset of the single lattice $\Lat$; the representation budget stays
$|\Plane|$ bits, shared, and each world adds one LUT and one threshold.
Consequently the diversity of languages costs nothing at the representation
level: the plane is the common earth, and the worlds are readings of it.
\end{proposition}

\begin{proof}
Proposition~\ref{prop:holographic}: the representation is common, the
interpretation is per-world, and the cost is a per-world LUT.
\end{proof}

\begin{remark}[Language and dialect are the same datatype]
Saussure refuses a principled line between language and dialect: the
difference is one of degree and of social-political history, not of
structure. The algebra agrees: a language and a dialect are both worlds; the
only structural differences are the size of the LUT, the generation of the
frame, and the threshold $\tau$. There is no type of world that is a
``language'' and another that is a ``dialect''.
\end{remark}

% ============================================================================
\section{Chapter II: Geographical diversity: its complexity}
\label{sec:ch2}
% ============================================================================

Saussure's chapter II complicates the picture: languages do not sit in clean
boxes. Dialects shade into one another, and the lines drawn between them---the
isoglosses---never coincide over any distance \cite[Part IV,
ch.~II]{harris1983}.

\begin{reading}[Continua and isoglosses]
\label{read:complexity}
The continuity of dialects is the graded overlap of interpretations: two
worlds whose LUTs agree on many tokens and whose frames are close read the
same region of the lattice similarly. An isogloss is a threshold on a
coordinate profile: the boundary where one soft key falls below $\tau$.
\end{reading}

\begin{definition}[Dialect continuum; isogloss]
\label{def:isogloss}
A \textbf{dialect continuum} is a family of worlds
$\World_1,\dots,\World_n$ such that adjacent worlds' interpretations overlap
on their shared nodes with BSS above threshold,
\[
  \BSS_{\mathcal{N}(\World_i)}\bigl(\mathcal{N}(\World_{i+1})\bigr)
  \;\ge\; \tau,
\]
for each $i$. An \textbf{isogloss} for coordinate $B_i$ of a frame is the
boundary of the set of communities whose soft key $w_i = \BSS_{B_i}$ exceeds
$\tau$: it is a level set of a graded quantity, not a cut.
\end{definition}

\begin{proposition}[Boundaries are thresholds, not cuts]
\label{prop:threshold}
For any two adjacent worlds in a continuum, the set of nodes on which they
agree is graded by the soft key: agreement is exact on the shared nodes
(Proposition~\ref{prop:sharing}) and partial elsewhere, and the soft key and
the spill measure how far apart the worlds are. Consequently there is no
point at which a language ``stops''; there are only thresholds below which
agreement is deemed insufficient. The complexity of geographical diversity is
the complexity of a graded boundary, not of a partition.
\end{proposition}

\begin{remark}[Why isoglosses never coincide]
Each isogloss is a level set of a different coordinate of a different frame;
different coordinates change at different rates under pivots, and different
frames are synchronized with different communities. There is no reason for
the level sets to coincide, and in the algebra they generically do not. This
is Saussure's observation that the bundle of isoglosses never forms a single
line, stated as the generic transversality of level sets.
\end{remark}

% ============================================================================
\section{Chapter III: Causes of geographical diversity}
\label{sec:ch3}
% ============================================================================

Saussure's chapter III asks why diversity exists and answers: time is the
primary cause. Geographical separation is not itself a cause; it is what
\emph{allows} divergence to accumulate without being reabsorbed
\cite[Part IV, ch.~III]{harris1983}.

\begin{reading}[Time as pivots; separation as unsynchronized generations]
\label{read:causes}
In the algebra, time is the sequence of pivots: every pivot increments a
generation and re-coordinates values. Two communities with a common ancestor
frame share the same langue exactly as long as they share the same pivots;
once their pivots stop being shared---once each community computes its own
extensions and rotations---their generations drift apart and their
transported coordinates diverge. Separation is not a force; it is the absence
of the bridge that would re-synchronize them.
\end{reading}

\begin{proposition}[Divergence is the product of transports]
\label{prop:divergence}
Let $\World_1$ and $\World_2$ share an ancestor frame $F(0)$, and suppose
that after time $0$ they undergo disjoint sequences of pivots with
change-of-basis matrices $\Change_1(1),\dots,\Change_1(t)$ and
$\Change_2(1),\dots,\Change_2(t)$. Then the value of a shared node $X$ at
time $t$ is
\[
  c_1 \;=\; c \,\Change_1(1)\cdots \Change_1(t),
  \qquad
  c_2 \;=\; c \,\Change_2(1)\cdots \Change_2(t),
\]
and the two communities diverge by the coordinate difference
$c_1 \xor c_2$. The divergence is zero exactly while the two sequences of
pivots coincide; from the first unsynchronized pivot onward, generically,
$c_1 \neq c_2$. Time, not space, is doing the work: the matrices are pure
generation increments.
\end{proposition}

\begin{proof}
Iterate the transport law of Part III \cite{part3}. Coordinate difference is
XOR in $\mathrm{GF}(2)^k$; the two sequences coincide if and only if every
pivot is shared.
\end{proof}

\begin{observation}[Substrata are old generations]
\label{obs:substratum}
When a wave of new speakers arrives and a community re-coordinates around
them, the older frame does not disappear: every generation remains readable
in the commit DAG, and every node keeps its old coordinates under the old
frame by exact transport. A \emph{substratum} is therefore not a mysterious
residue; it is the set of old generations still referenced in the DAG. The
algebra keeps substrata as first-class objects: any node can be read against
the frame of any earlier generation, with the spill measuring what that old
lens cannot see \cite{state2026}.
\end{observation}

\begin{remark}[Saussure's refusal of internal causes]
Saussure denies that geographical diversity has an internal, linguistic
cause. The algebra agrees: the pivot is occasioned by residue (Part III), and
residue arrives from parole---from outside the system. Geography merely
decides which pivots are shared and which are not. The cause of diversity is
the same residual that Part III located, allowed to accumulate in separated
worlds.
\end{remark}

% ============================================================================
\section{Chapter IV: Propagation of linguistic waves}
\label{sec:ch4}
% ============================================================================

Saussure's chapter IV describes how innovations spread: a linguistic wave
travels across a territory through communication, and it may be stopped by
barriers \cite[Part IV, ch.~IV]{harris1983}.

\begin{reading}[The wave is the bridge]
\label{read:wave}
A linguistic wave is the bridge of Definition~\ref{def:bridge}. An
innovation---a node, or a pivot with its new opposition---crosses from one
world to another by being materialized in the source, transcribed into shared
encodings, and ingested in the target. Barriers are LUT separation (no shared
codebook across which to transcribe) and threshold rejection (the imported
node's soft key falls below $\tau$ in the target frame).
\end{reading}

\begin{proposition}[Waves propagate without loss on builders]
\label{prop:wave}
Let $\Bridge = \Bridge_{\World_1\to\World_2}$ be the bridge. Then
$\Bridge$ is deterministic, idempotent on the builders of the source nodes,
and monotone; hence an innovation propagates exactly to the extent that its
builders survive materialization, and it can only gain collateral atoms from
collisions, never lose a builder. A wave is therefore a monotone, convergent
spread: it does not weaken with distance, it is \emph{blocked}---by a missing
codebook or a threshold---or it arrives whole.
\end{proposition}

\begin{proof}
The bridge laws of Proposition~\ref{prop:bridge-laws}, restated
geographically.
\end{proof}

\begin{observation}[Loans and inheritances are separated by provenance]
\label{obs:loans}
A node inherited from a common ancestor has a provenance path inside one DAG
lineage: its recipe names only ancestors of that lineage. A node borrowed
across worlds has a provenance path that \emph{crosses} the bridge: its
recipe names a source node in $\World_1$ and a re-ingestion in $\World_2$.
The algebra therefore separates loans from inheritances by provenance alone,
without any external etymological dictionary. A loan is a view whose sources
live in two worlds; an inheritance is a view whose sources live in one
lineage.
\end{observation}

\begin{remark}[The dispatcher is the speaker]
Part II \cite{part2} introduced the dispatcher
$\Disp : \Lat \times \mathrm{State} \to \Delta(\mathrm{Hosts})$ as the
mechanism that selects which host actualizes the shared langue. In the
geographical reading, the dispatcher is the individual speaker who chooses,
at each utterance, which world to speak: the learned policy over hosts is the
executive side of the linguistic community's geography. Langue never speaks;
worlds never speak; the dispatcher speaks, and its choices are the daily
traffic in which waves travel.
\end{remark}

% ============================================================================
\section{The dictionary of Part Four}
\label{sec:dictionary}
% ============================================================================

\begin{table}[htbp]
\centering
\renewcommand{\arraystretch}{1.2}
\begin{tabular}{@{}>{\raggedright\arraybackslash}p{3.6cm}
                    >{\raggedright\arraybackslash}p{4.2cm}
                    >{\raggedright\arraybackslash}p{6.8cm}@{}}
\toprule
Saussure (Part Four) & In the \textit{Course} & In HLLSet Algebra \\
\midrule
diversity of languages & languages differ in space & the family of worlds
$\Wset$ over one plane: Reading~\ref{read:diversity},
Proposition~\ref{prop:one-lattice} \\
language vs dialect & difference of degree & one datatype (world); only LUT
size, generation, and threshold differ \\
geographical complexity & dialect continua, isoglosses & graded overlap and
threshold level sets: Definition~\ref{def:isogloss},
Proposition~\ref{prop:threshold} \\
cause of diversity & time, not space & unsynchronized pivots; divergence is
the product of transports: Proposition~\ref{prop:divergence} \\
substratum & an older layer under the present & old generations readable in
the DAG: Observation~\ref{obs:substratum} \\
linguistic wave & innovation spreading by communication & the bridge
$\Bridge$: Definition~\ref{def:bridge}, Proposition~\ref{prop:wave} \\
barrier to a wave & geographical/social obstacle & missing codebook or
threshold rejection: Reading~\ref{read:wave} \\
borrowing vs inheritance & loans and native forms & provenance crossing two
worlds vs one lineage: Observation~\ref{obs:loans} \\
the speaker & who actualizes langue & the dispatcher $\Disp$ over hosts:
Part II \cite{part2}, Remark on the dispatcher \\
\bottomrule
\end{tabular}
\caption{The complete dictionary of Part Four of the \textit{Course} into
HLLSet Algebra.}
\label{tab:dictionary}
\end{table}

% ============================================================================
\FloatBarrier
\section{What Part IV adds to Parts I--III}
\label{sec:adds}
% ============================================================================

\begin{enumerate}[leftmargin=*,nosep]
  \item \textbf{The world layer becomes a linguistic object.} The EWM
    multi-world machinery \cite{companion2026} was built for isolation and
    bridges; Part IV reads it as Saussure's geography: worlds are languages,
    LUTs are lexicons, thresholds are community boundaries.
  \item \textbf{Diversity is costless at the representation level.}
    Proposition~\ref{prop:one-lattice} is the exact form of the
    structuralist claim that all languages are one system seen differently.
  \item \textbf{The cause of diversity is time, computed.}
    Proposition~\ref{prop:divergence} turns Saussure's ``time is the primary
    cause'' into a product of change-of-basis matrices.
  \item \textbf{Loans and inheritances are decidable.}
    Observation~\ref{obs:loans} separates borrowed and inherited nodes by
    provenance, without a dictionary.
  \item \textbf{The dispatcher finds its place.} The executive question of
    Part II becomes the geography of the speaker.
\end{enumerate}

% ============================================================================
\section{Conclusion and outlook to Part V}
\label{sec:conclusion}
% ============================================================================

Part Four of the \textit{Course} has been re-founded as the multi-world
theory of one lattice. Diversity is a family of worlds over one plane;
complexity is graded overlap with threshold boundaries; the cause of
diversity is time, made exact as unsynchronized pivots; waves are bridges;
substrata are old generations; loans are crossing provenances. Geographical
linguistics, read in this direction, is no longer a catalogue of dialects: it
is the theory of which worlds share which frames, and how innovations cross
between them.

What Part Four leaves to Part Five is the direction of reading itself. We
have read the DAG forward: pivots, transports, waves. Part Five of the
\textit{Course}, Questions of Retrospective Linguistics, reads it backward:
the two perspectives of diachronic linguistics, the search for the earliest
language and the prototype, reconstructions, the contribution of linguistic
evidence to anthropology and prehistory, and the distinction of language
families from linguistic types. In the algebra these become time travel
through old frames: reading the present against the first recorded basis,
with the spill as the error of the journey; reconstruction as a view whose
provenance points into the past; and the family/type distinction as ancestry
in the DAG versus similarity in the current fiber. Part V will close the
series, and with it the circle: the \textit{Course} interpreted by the
algebra, and the algebra interpreted by the \textit{Course}.

\section*{Acknowledgments}
This paper was created in collaboration with, and with the assistance of, an
AI system---DeepSeek (deepseek-v4-pro). The world-layer reading of
geographical diversity, the transport-product reading of divergence, the
provenance reading of loans, and several of the formulations emerged in
dialogue with the assistant, whose drafts the authors revised and take full
responsibility for in their final form.

\begin{thebibliography}{99}

\bibitem{saussure1916} F.~de Saussure,
\textit{Cours de linguistique g\'en\'erale}, Payot, Paris, 1916.
(Edited posthumously by C.~Bally and A.~Sechehaye, with A.~Riedlinger.)

\bibitem{harris1983} F.~de Saussure,
\textit{Course in General Linguistics}, translated and annotated by
Roy Harris, Duckworth, London, 1983.

\bibitem{baskin1959} F.~de Saussure,
\textit{Course in General Linguistics}, translated by Wade Baskin,
Philosophical Library, New York, 1959.

\bibitem{saussure1977} F.~de Sossyur,
\textit{Trudy po yazykoznaniyu} [Works in Linguistics], edited by
A.~A.~Kholodovich, translation of the \textit{Cours} by A.~M.~Sukhotin,
Progress, Moscow, 1977.
(Authors' working copy; used for verification of the readings.)

\bibitem{part1} A.~Mylnikov, A.~Mylnikov Jr., and D.~Kumar,
\textit{Saussure in HLLSet Algebra, Part I: General Principles},
SAUSSURE PUB series, 2026.

\bibitem{part2} A.~Mylnikov, A.~Mylnikov Jr., and D.~Kumar,
\textit{Saussure in HLLSet Algebra, Part II: Synchronic Linguistics},
SAUSSURE PUB series, 2026.

\bibitem{part3} A.~Mylnikov, A.~Mylnikov Jr., and D.~Kumar,
\textit{Saussure in HLLSet Algebra, Part III: Diachronic Linguistics},
SAUSSURE PUB series, 2026.

\bibitem{ewm2026} A.~Mylnikov, A.~Mylnikov Jr., and D.~Kumar,
\textit{Emerging World Models: A Structural Framework for Self-Organizing
World Representations}, EWM PUB series, 2026.

\bibitem{companion2026} A.~Mylnikov, A.~Mylnikov Jr., and D.~Kumar,
\textit{The EWM Lattice and Its Worlds}, EWM PUB series, 2026.

\bibitem{state2026} A.~Mylnikov, A.~Mylnikov Jr., and D.~Kumar,
\textit{Transformations and Trajectories of a Content-Addressed State
Machine}, EWM PUB series, 2026.

\bibitem{attention2026} A.~Mylnikov, A.~Mylnikov Jr., and D.~Kumar,
\textit{Attention on Lattice Traversal Is All You Need}, EWM PUB series, 2026.

\bibitem{langue2026} A.~Mylnikov, A.~Mylnikov Jr., and D.~Kumar,
\textit{Langue as a Lattice: Closing the Circle from Statistical to
Structural Linguistics through Content-Addressed Algebras}, EWM PUB series,
2026.

\bibitem{birkhoff1940} G.~Birkhoff,
\textit{Lattice Theory}, American Mathematical Society, 1940.

\bibitem{maclane1971} S.~Mac Lane,
\textit{Categories for the Working Mathematician}, Springer, 1971.

\bibitem{hjelmslev1961} L.~Hjelmslev,
\textit{Prolegomena to a Theory of Language}, University of Wisconsin Press,
1961.

\bibitem{harris1954} Z.~S.~Harris,
\textit{Distributional Structure}, Word, 10(2--3):146--162, 1954.

\bibitem{firth1957} J.~R.~Firth,
\textit{A Synopsis of Linguistic Theory, 1930--1955}, Studies in Linguistic
Analysis, 1957.

\bibitem{vaswani2017} A.~Vaswani, N.~Shazeer, N.~Parmar, J.~Uszkoreit,
L.~Jones, A.~N.~Gomez, \L.~Kaiser, and I.~Polosukhin,
\textit{Attention Is All You Need}, NeurIPS, 2017.

\end{thebibliography}

\end{document}
